\chapter{Dividends, Borrow and Forwards}\label{ch:dv:dividends-borrow-and-forwards}

A share at 50 is on every broker's hard-to-borrow list. Its one-month options
look broken: at the strike of 50 the put trades at 4.04 and the call at 2.78,
and a screen that computes implied volatilities from the spot and the interest
rate shows the put at 72 volatility and the call at 47. A new trader sees 25
points of arbitrage: sell the dear put, buy the cheap call, hedge. There is
none. Put--call parity (One Quant Book 1, chapter 25) says the gap between the
call and the put is a forward, and the forward the options are pricing is
48.74, not the 50.16 that the spot and the rate suggest: the market is charging
35\% a year to borrow the share, and anyone who wants to be short it through
the options pays that rate either way. This chapter is about the forward: what
goes into it for a share (dividends and the cost of borrowing), how dividends
enter an option model, how the chain itself reveals the forward, and who carries
the risk that dividends change.

\section{Forwards with discrete dividends and borrow}

A share pays dividends on known ex-dividend dates (One Quant Book 1, chapter 8),
and a short seller pays a borrow fee to the lender (One Quant Book 1, chapter 6).
Both reduce the forward below the financed spot: the buyer of a forward does not
receive the dividends paid before delivery, and the seller who holds the share as
a hedge can lend it out and earn the fee.

\begin{proposition}[The equity forward]\label{prop:dv:dividends-borrow-and-forwards:fwd}
With cash dividends $D_i$ paid at $t_i$, a flat financing rate $r$ and a borrow fee
$\ell$ (continuously compounded, paid on the value lent), the forward for delivery
at $T$ is
\[
F_{0,T}=\Bigl(S_0-\sum_{t_i\le T}D_iP(0,t_i)\Bigr)\frac{e^{-\ell T}}{P(0,T)} .
\]
\end{proposition}
\begin{proof}
Buy the share for $S_0$, borrow the present value of the dividends against them,
lend the share out and reinvest the fee, sell the forward. At $T$ the position is
flat and the dividends have repaid their loan: the cost of the share net of what it
earns must equal the present value of the delivery price.
\end{proof}

\begin{definition}[Proportional dividend]\label{def:dv:dividends-borrow-and-forwards:prop}
A \emph{proportional dividend}\index{proportional dividend} is modelled as a fixed
fraction $y_j$ of the share price at the ex-date rather than a fixed amount: the
share drops by $y_jS_{t_j-}$, and the forward is multiplied by $(1-y_j)$ for every
ex-date before $T$. A continuous dividend yield $q$ is its limit when payments are
spread evenly over time.
\end{definition}

The forward curve of a share is a sawtooth
(\cref{fig:dv:dividends-borrow-and-forwards:curve}): it grows at the financing
rate less the borrow between ex-dates and drops by each dividend. A continuous yield
with the same two-year forward is a straight line through it, wrong at every
intermediate date by up to half a quarter's dividend. For options whose expiry
straddles an ex-date the difference is the whole dividend.

\begin{omfigure}
\begin{tikzpicture}
\begin{axis}[width=10.5cm,height=4.8cm,
  xlabel={delivery date (years)},ylabel={forward},
  tick label style={font=\small},label style={font=\small},
  xmin=0,xmax=2,ymin=99.3,ymax=100.9,grid=major,grid style={gray!25},
  yticklabel style={/pgf/number format/.cd,fixed,precision=1},
  legend columns=2,legend style={font=\footnotesize,at={(0.5,-0.3)},anchor=north,draw=none,fill=none,/tikz/every even column/.append style={column sep=8pt}},
  table/col sep=comma]
\addplot[omDef,very thick,const plot mark left] table[x=t,y=discrete]{figdata/derivatives/05-dividends-borrow-and-forwards/forward_curve.csv};
\addplot[omThm,thick,dashed] table[x=t,y=smooth]{figdata/derivatives/05-dividends-borrow-and-forwards/forward_curve.csv};
\legend{quarterly cash dividends of 0.60,continuous yield with the same 2-year forward}
\end{axis}
\end{tikzpicture}

\omcaption{Forward curve of a share at 100 paying 0.60 a quarter, with a 3\%
financing rate and a 0.5\% borrow fee. Between ex-dates the forward grows at 2.5\%
a year; on each ex-date it drops by the dividend. Illustrative
parameters.}\label{fig:dv:dividends-borrow-and-forwards:curve}
\end{omfigure}

\section{Dividends in the option model}

For a European option, only the forward and the distribution around it matter; the
Black formula of chapter 3 prices it on the forward. The choice of dividend model
decides what the quoted volatility is the volatility of, and how the forward moves
when the spot moves.

\begin{definition}[Escrowed dividend model]\label{def:dv:dividends-borrow-and-forwards:escrowed}
In the \emph{escrowed dividend model}\index{escrowed dividend model} the share is
split into the present value of the dividends to be paid before expiry, treated as
riskless, and the rest, $S^*_t=S_t-\sum_{t<t_i\le T}D_iP(t,t_i)$, which follows a
geometric Brownian motion. European options are priced by Black--Scholes with
$S^*_0$ in place of $S_0$.
\end{definition}

Three models are in use. In the escrowed model the volatility is that of $S^*$.
In the proportional model it is that of the whole share, which stays lognormal
through its proportional drops. In the \emph{spot model} the whole share is
lognormal between ex-dates and drops by the cash amount on each: the most
realistic, and the only one with no closed form. For a European option with the
same forward, the escrowed and proportional models give the same price at the same
volatility, since both make the forward lognormal with that volatility; the spot
model gives a higher price, because before the ex-date the part of the share that
will be paid out as a dividend also moves.

\begin{proposition}[The volatility of the escrowed share]\label{prop:dv:dividends-borrow-and-forwards:adj}
If the share has volatility $\sigma$ in the spot model and pays one dividend of
present value $\bar D$ at $t_1<T$, the escrowed model reproduces its at-the-money
price to first order with the effective volatility
\[
\sigma_{\mathrm{eff}}^2\,T=\sigma^2\Bigl(\frac{S_0}{S_0-\bar D}\Bigr)^2t_1+\sigma^2(T-t_1).
\]
\end{proposition}
\begin{proof}
Before $t_1$ the escrowed share $S^*=S-\bar D e^{rt}$ moves by the whole share's
moves, $dS^*=dS=\sigma S\,dW$, so its relative volatility is $\sigma S/S^*$; after
$t_1$ the two coincide. Adding the variances of the two periods with $S/S^*$ frozen
at its initial value gives the formula.
\end{proof}

\begin{example}[One dividend of 4]\label{ex:dv:dividends-borrow-and-forwards:adj}
Share 100, one-year options, $r=3\%$, a dividend of 4 in six months, spot-model
volatility 25\%. The escrowed and proportional models at 25\% imply a flat 25.00;
the spot model's prices imply 25.57 at the 70 strike and 25.49 at the 130; the
effective volatility of the proposition is 25.52
(\cref{fig:dv:dividends-borrow-and-forwards:models}). Half a volatility point is
0.2 of premium at the money: a desk that mixes models across systems books it as
profit or loss.
\end{example}

\begin{omfigure}
\begin{tikzpicture}
\begin{axis}[width=10.5cm,height=4.8cm,
  xlabel={strike},ylabel={implied volatility (\%)},
  tick label style={font=\small},label style={font=\small},
  xmin=70,xmax=130,ymin=24.9,ymax=25.7,grid=major,grid style={gray!25},
  yticklabel style={/pgf/number format/.cd,fixed,precision=1},
  legend columns=2,legend style={font=\footnotesize,at={(0.5,-0.3)},anchor=north,draw=none,fill=none,/tikz/every even column/.append style={column sep=8pt}},
  table/col sep=comma]
\addplot[omDef,very thick] table[x=strike,y=escrowed]{figdata/derivatives/05-dividends-borrow-and-forwards/model_smiles.csv};
\addplot[omThm,very thick,mark=*,mark size=1.5pt] table[x=strike,y=spot]{figdata/derivatives/05-dividends-borrow-and-forwards/model_smiles.csv};
\addplot[omProp,thick,dashed] table[x=strike,y=adjusted]{figdata/derivatives/05-dividends-borrow-and-forwards/model_smiles.csv};
\legend{escrowed or proportional at 25\%,spot model at 25\%,escrowed at the effective volatility}
\end{axis}
\end{tikzpicture}

\omcaption{One-year calls on a share at 100 with a dividend of 4 in six months,
priced by three dividend models with the same 25\% volatility and read back as
Black implied volatilities on the common forward. The spot model is half a point
above the others and slightly skewed; the effective volatility of
\cref{prop:dv:dividends-borrow-and-forwards:adj} matches it at the money. Data:
the tutorial.}\label{fig:dv:dividends-borrow-and-forwards:models}
\end{omfigure}

\begin{definition}[Mixed dividend model]\label{def:dv:dividends-borrow-and-forwards:mixed}
A \emph{mixed dividend model}\index{mixed dividend model} treats near dividends as
cash amounts, announced or forecast with confidence, and dividends beyond a horizon
(one to two years) as proportional to the share price, blending the two over a
transition period.
\end{definition}

The blend is about dynamics. With cash dividends the forward moves one for one
with the spot, $\partial F/\partial S=1/P(0,T)$; with \omterm{def:dv:dividends-borrow-and-forwards:prop}{proportional dividends} it
moves in proportion, $\partial F/\partial S=F/S$. For a five-year option on a share
yielding 4\% the two deltas differ by almost a fifth, and so does the hedge. A company
whose share halves rarely keeps paying the same cash dividend, which is why far
dividends are modelled as proportional; near dividends are already declared.

\section{Implying forwards, dividends and borrow from parity}

For European options put--call parity, $C-P=P(0,T)(F-K)$, holds at every strike.
It is linear in the strike, with slope $-P(0,T)$ and intercept $P(0,T)F$, so a
regression across strikes reads both off the chain without any model.

\begin{method}[Reading the forward curve off a chain]\label{met:dv:dividends-borrow-and-forwards:read}
For each expiry:
\begin{enumerate}
\item regress $C_K-P_K$ on $K$ over the strikes near the money: the slope is
$-P(0,T)$, the intercept $P(0,T)F$;
\item the implied carry $S-FP(0,T)$ is the present value of dividends plus borrow
cost to $T$; its increments between expiries give the dividend term structure;
\item with a dividend forecast, what remains of the carry is the borrow, the next
definition; with the borrow known (general collateral), what remains is the market's
dividend forecast.
\end{enumerate}
For American single-stock options parity is only a pair of inequalities; use
European index options, or remove the early-exercise premium first (chapter 6).
\end{method}

\begin{definition}[Implied dividend and implied borrow rate]\label{def:dv:dividends-borrow-and-forwards:implied}
The \emph{implied dividend}\index{implied dividend} of an expiry is the present value
of dividends to that expiry that makes the forward read from options, together with
the known borrow cost, consistent with the spot. The \emph{implied borrow
rate}\index{implied borrow rate} is the borrow fee $\ell$ that makes the forward
consistent with the spot and a dividend forecast:
$\ell=-\ln\bigl(FP(0,T)/(S-\bar D)\bigr)/T$.
\end{definition}

\begin{example}[The hard-to-borrow share]\label{ex:dv:dividends-borrow-and-forwards:htb}
The one-month chain of the opening, five strikes from 45 to 55 quoted to the cent,
regressed as in \cref{met:dv:dividends-borrow-and-forwards:read}, gives a forward of
48.74 and a discount factor of 0.99680 (a rate of 3.90\%, against 4\% true: the
cents of rounding). With no dividend due, the carry of 1.42 is all borrow: 35.0\% a
year. Priced from the naive forward $50e^{0.04\times30/365}=50.16$ at the chain's 60\%
volatility, the call would be worth 3.50 instead of 2.78, and the put 3.34 instead of
4.04. The ``skew'' of 72 against 47 is the wrong forward, not the market's view of
volatility.
\end{example}

Selling the put, buying the call and shorting the share (the reversal of One Quant
Book 1, chapter 25) appears to earn 1.42 in a month; it requires borrowing the share
for a month at 35\%, which costs 1.44. Options on hard-to-borrow shares are where
parity fails most visibly when measured from the spot, and academic studies find that
the size of the apparent violation tracks the stock-lending market's fee
(Ofek, Richardson and Whitelaw).

\begin{omfigure}
\begin{tikzpicture}
\begin{axis}[width=10.5cm,height=4.8cm,ybar,bar width=10pt,
  xlabel={expiry (years)},ylabel={carry per quarter},
  tick label style={font=\small},label style={font=\small},
  xmin=0.1,xmax=2.15,ymin=0.5,ymax=0.9,ymajorgrids,grid style={gray!25},
  xtick={0.25,0.5,0.75,1,1.25,1.5,1.75,2},
  legend columns=2,legend style={font=\footnotesize,at={(0.5,-0.36)},anchor=north,draw=none,fill=none,/tikz/every even column/.append style={column sep=8pt}},
  table/col sep=comma]
\addplot[fill=omDef!55,draw=omDef] table[x=t,y=implied]{figdata/derivatives/05-dividends-borrow-and-forwards/carry_steps.csv};
\addplot[fill=omThm!45,draw=omThm] table[x=t,y=true]{figdata/derivatives/05-dividends-borrow-and-forwards/carry_steps.csv};
\legend{read from parity (quotes $\pm$0.02),true: PV of the dividend plus borrow}
\end{axis}
\end{tikzpicture}

\omcaption{Present value of the carry (dividend plus borrow cost) between consecutive
quarterly expiries, recovered by the parity regression from a synthetic chain whose
mid prices carry up to two cents of noise at nine strikes. The cumulative forward is
recovered to about a cent; each quarter's increment only to about 0.08. Data: the
tutorial.}\label{fig:dv:dividends-borrow-and-forwards:carry}
\end{omfigure}

\section{Dividend risk and dividend swaps}

\begin{definition}[Dividend risk and dividend swap]\label{def:dv:dividends-borrow-and-forwards:risk}
\emph{Dividend risk}\index{dividend risk} is the sensitivity of a position to a change
in the dividends expected before its maturity: a forward, a long call or a short put
loses when expected dividends rise. A \emph{dividend swap}\index{dividend swap} is an
over-the-counter contract that exchanges, at maturity, the dividends actually paid on
a share or index over a period for a fixed amount agreed at inception; the listed
version is the dividend future of One Quant Book 1, chapter 22.
\end{definition}

\omterm{def:dv:dividends-borrow-and-forwards:risk}{Dividend risk} is small for a one-month option and dominant for a five-year one. It
accumulates on the books of structured-product issuers: a bank that sells notes
paying a function of an index hedges by holding the index, collects its dividends, and
has priced those dividends into the notes; if dividends fall, the hedge earns less than
was promised. Practitioners describe such issuers as structurally long dividends and
as sellers of \omterm{def:dv:dividends-borrow-and-forwards:risk}{dividend swaps} and futures to lay the risk off (chapter 18). The buyers
are investors who want dividend exposure without equity exposure.

\begin{dated}{2026-09}{Index dividend futures and 2020}\label{dat:dv:dividends-borrow-and-forwards:divfut}
The euro-area benchmark index dividend futures listed on Eurex have a contract value
of EUR~100 per index point and a tick of 0.1 point (EUR~10), settle in cash on the
cumulative gross dividends paid by the index members over the contract year, and list
annual maturities up to about ten years out. On 27 March 2020 the European Central Bank
recommended that the banks it supervises pay no dividends until at least 1 October
2020, later extended to January 2021. A practitioner account of April 2020 reported
2020 index dividends trading at 55\% of what models had predicted and 2021 dividends
at 70\%, as issuers sold dividends into a falling market.
\end{dated}

\begin{omfigure}
\begin{tikzpicture}[font=\small,
  bx/.style={draw=omDef,thick,rounded corners=2pt,fill=omDef!6,align=center,text width=2.6cm,inner sep=4pt}]
\node[bx] (inv) at (0,0) {investors in\\structured notes};
\node[bx] (iss) at (5.6,0) {issuer: holds the\\index as hedge,\\long dividends};
\node[bx] (buy) at (11.2,0) {dividend buyers:\\funds, pension\\funds};
\node[bx] (co) at (5.6,-2.6) {index members\\pay the dividends};
\draw[-{Stealth},thick] (inv) -- (iss);
\draw[-{Stealth},thick] (iss) -- (buy);
\node[font=\footnotesize,align=center,anchor=south] at (2.8,0.1) {buy notes priced\\with dividends};
\node[font=\footnotesize,align=center,anchor=south] at (8.4,0.1) {sell dividend\\swaps, futures};
\draw[-{Stealth},thick,omThm] (co) -- node[right,font=\footnotesize]{dividends} (iss);
\end{tikzpicture}

\omcaption{Where \omterm{def:dv:dividends-borrow-and-forwards:risk}{dividend risk} goes. The issuer of equity-linked notes hedges with the
underlying, is paid its dividends and has promised their expected value to the note
holders; it lays the risk off by selling dividends forward to investors who want
them.}\label{fig:dv:dividends-borrow-and-forwards:flows}
\end{omfigure}

\section{Tutorial: a forward curve from a chain}\label{tut:dv:dividends-borrow-and-forwards:tut}

\begin{tutorial}
\textbf{Goal.} Read forwards, discount factors, carry and borrow off a synthetic chain,
and compare the three dividend models on one option. \textbf{End state:}
\cref{fig:dv:dividends-borrow-and-forwards:models,fig:dv:dividends-borrow-and-forwards:carry}
and the numbers of \cref{ex:dv:dividends-borrow-and-forwards:htb}.
\begin{tutsteps}
\item \textbf{Parity by regression}, then carry, borrow and the dividend strip:
\omcode{code/firm/divfwd/firm_divfwd.py}{34}{58}{Forward and discount factor from a chain; implied carry, borrow and dividend strip.}\label{lst:dv:dividends-borrow-and-forwards:parity}
\item \textbf{The spot model} by Gauss--Hermite quadrature over the share just before
the ex-date, and the effective volatility of the escrowed share:
\omcode{code/derivatives/05-dividends-borrow-and-forwards/python/dv_dividends.py}{35}{53}{The spot model, and the escrowed model's effective volatility.}\label{lst:dv:dividends-borrow-and-forwards:spot}
\item \textbf{Run} \texttt{dv\_dividends.curve\_from\_chain()},
\texttt{htb\_results()} and \texttt{fig\_dividends.py}.
\end{tutsteps}
\textbf{What to change next.} Regress on the five strikes nearest the forward only, and
on all strikes weighted by the inverse of the quote width, and compare the noise of the
quarterly carry; then give the share two dividends before expiry and extend the
effective volatility to both.
\end{tutorial}

\section{Build: the equity forward curve}\label{bld:dv:dividends-borrow-and-forwards:divfwd}

\begin{build}
\bfield{Purpose} Every equity option of the miniature firm is priced on a forward that
comes from this curve; in the pricing library of chapter 28 it is what
\texttt{MarketData.forward} computes.
\bfield{Interface} \texttt{ForwardCurve(spot, rate, cash, proportional,
borrow).forward(t)}; \texttt{implied\_forward(strikes, calls, puts) -> (F, df)};
\texttt{implied\_carry}; \texttt{implied\_borrow(spot, fwd, df, t, pv\_dividends)};
\texttt{strip\_dividends(spot, expiries, forwards, dfs)}.
\bfield{Rules} Dividends strictly after today and on or before the delivery date count;
\omterm{def:dv:dividends-borrow-and-forwards:prop}{proportional dividends} multiply the forward; the regression uses European quotes (or
de-Americanised ones).
\bfield{Acceptance tests} \texttt{code/firm/divfwd/tests/}: the forward replicates by the
cash-and-carry argument; the regression recovers a known forward and discount factor
exactly from noiseless quotes; the implied borrow of a clean chain equals the fee used
to build it; strips sum to the total carry.
\bfield{Stretch} Fit a smooth borrow curve and a dividend curve jointly across expiries
with a penalty on roughness, and flag expiries whose parity residuals are outliers.
\end{build}

\begin{omsources}
\item E.~G.~Haug, J.~Haug and A.~Lewis, ``Back to basics: a new approach to the
discrete dividend problem'', \emph{Wilmott} (September 2003) 37--47.
\item E.~Ofek, M.~Richardson and R.~F.~Whitelaw, ``Limited arbitrage and short sales
restrictions: evidence from the options markets'', \emph{Journal of Financial
Economics} 74 (2004) 305--342.
\item European Central Bank, Recommendation ECB/2020/19 of 27 March 2020 on dividend
distributions during the COVID-19 pandemic.
\item Eurex, EURO STOXX~50 Index Dividend Futures, contract specifications.
\item E.~Barthe, ``Hedging structured products during the corona crisis: where did it
go wrong?'', Structured Retail Products, 30 April 2020.
\end{omsources}

\section{Exercises}

\begin{exercise}[$\star$]\label{exo:dv:dividends-borrow-and-forwards:1}
A share at 100 pays dividends of 1 in three and in nine months; $r=3\%$, no borrow
fee. Give the one-year forward.
\end{exercise}

\begin{exercise}[$\star$]\label{exo:dv:dividends-borrow-and-forwards:2}
At one expiry, $C-P$ is 7.02 at the strike 95 and $-2.83$ at the strike 105. Give the
discount factor and the forward.
\end{exercise}

\begin{exercise}[$\star$]\label{exo:dv:dividends-borrow-and-forwards:3}
Check the effective volatility 25.52\% of \cref{ex:dv:dividends-borrow-and-forwards:adj}.
\end{exercise}

\begin{exercise}[$\star\star$]\label{exo:dv:dividends-borrow-and-forwards:4}
Why does put--call parity give only bounds for American single-stock options, and in
which direction is the error when the European relation is used near an ex-date?
\end{exercise}

\begin{exercise}[$\star\star$]\label{exo:dv:dividends-borrow-and-forwards:5}
A share at 100 pays a dividend worth 4\% of its price in six months; $r=3\%$, one year.
Give $\partial F/\partial S$ when the dividend is modelled as the cash amount 4 and when
it is modelled as proportional.
\end{exercise}

\begin{exercise}[$\star\star$]\label{exo:dv:dividends-borrow-and-forwards:6}
A three-year at-the-money call on the share of \cref{fig:dv:dividends-borrow-and-forwards:curve}
(quarterly dividends of 0.60, $r=3\%$, borrow 0.5\%, volatility 20\%) is repriced after a
10\% cut in every future dividend. By how much does its value change?
\end{exercise}

\begin{exercise}[$\star\star\star$]\label{exo:dv:dividends-borrow-and-forwards:7}
\emph{Coding.} From \texttt{synthetic\_chain()}, compute the implied forward and
discount factor of the one-year expiry and the implied borrow given the known dividends,
and compare with the 0.5\% used to build the chain.
\end{exercise}

\begin{exercise}[$\star\star\star$]\label{exo:dv:dividends-borrow-and-forwards:8}
\emph{Find the flaw.} ``On this hard-to-borrow share the one-month 50 puts trade at 72
volatility and the calls at 47. Selling puts and buying calls, delta-hedged, captures 25
volatility points.''
\end{exercise}

\section{Problem: The Hard-to-Borrow Share}

\begin{problem}[{Weekend problem --- reading the borrow off the options}]\label{pb:dv:dividends-borrow-and-forwards:1}
The share of the opening trades at 50 and pays no dividend in the next month; the
financing rate is 4\%. Its one-month European options are quoted (mid): 45 call 5.41,
put 1.68; 47.5 call 3.94, put 2.71; 50 call 2.78, put 4.04; 52.5 call 1.90, put 5.65; 55
call 1.26, put 7.50. The chain's at-the-money volatility is 60\%.

\medskip\noindent\textbf{Part I --- The forward.}
\begin{enumerate}
\item Compute $C-P$ at each strike.
\item Regress it on the strike: give the discount factor and the forward.
\item Give the implied rate, and explain the gap to 4\%.
\item Give the implied carry.
\item Give the \omterm{def:dv:dividends-borrow-and-forwards:implied}{implied borrow rate}.
\end{enumerate}

\medskip\noindent\textbf{Part II --- The wrong forward.}
\begin{enumerate}[resume]
\item Give the naive forward from the spot and the rate.
\item Price the 50 call and put at 60\% on the naive forward.
\item By how much are they mispriced?
\item Give their implied volatilities computed on the naive forward.
\item Why do the two implied volatilities differ, and which forward makes them equal?
\end{enumerate}

\medskip\noindent\textbf{Part III --- The apparent arbitrage.}
\begin{enumerate}[resume]
\item What does the reversal (short share, short put, long call) appear to earn?
\item What does it cost to borrow the share for the month?
\item Is there an arbitrage once borrow is counted?
\item What can a holder of the share do with the options instead of lending it, and what
does that earn?
\item How would a desk that is not able to borrow the share express a view that it will fall?
\end{enumerate}

\medskip\noindent\textbf{Part IV --- Judgement.}
\begin{enumerate}[resume]
\item What happens to the options if the borrow fee falls to 5\% overnight?
\item Why are listed options on such shares American, and what does that change?
\item What should an options system use as the forward for this share?
\item State the \emph{named result}: the \omterm{def:dv:dividends-borrow-and-forwards:implied}{implied borrow rate} and the naive pricing error on
the 50 call.
\item In one sentence: what does put--call parity measure on a hard-to-borrow share?
\end{enumerate}
\end{problem}

\section{Interview questions}

\begin{interviewq}[$\star$ \iqroles{trader, researcher}]\label{iq:dv:dividends-borrow-and-forwards:1}
How do you get the forward of a share from its option prices, without a model?
\end{interviewq}

\begin{interviewq}[$\star$ \iqroles{trader}]\label{iq:dv:dividends-borrow-and-forwards:2}
The at-the-money put of a share trades well above the call at the same strike. Give two
reasons.
\end{interviewq}

\begin{interviewq}[$\star\star$ \iqroles{researcher, bank}]\label{iq:dv:dividends-borrow-and-forwards:3}
Cash or \omterm{def:dv:dividends-borrow-and-forwards:prop}{proportional dividends}: which do you use for a one-month option and for a
five-year one, and what changes in the hedge?
\end{interviewq}

\begin{interviewq}[$\star\star$ \iqroles{trader, risk}]\label{iq:dv:dividends-borrow-and-forwards:4}
What is \omterm{def:dv:dividends-borrow-and-forwards:risk}{dividend risk}, who is naturally long it, and what happened to it in 2020?
\end{interviewq}

\begin{interviewq}[$\star\star$ \iqroles{developer, researcher}]\label{iq:dv:dividends-borrow-and-forwards:5}
How do you put a discrete cash dividend into a tree pricer for American options, and
what goes wrong with the obvious approaches?
\end{interviewq}

\begin{interviewq}[$\star\star\star$ \iqroles{researcher}]\label{iq:dv:dividends-borrow-and-forwards:6}
The same volatility of 25\% is fed to the escrowed model in one system and the spot model
in another. Which gives the higher call price, by roughly how much, and why?
\end{interviewq}
